% Source: https://artofproblemsolving.com/community/c3h1266896_number_theory_and_perfect_squares
% Stable identity: aops-topic-1266896-q01
Note that $n(n+2)=(n+1)^2-1$. Any common divisor $d$ of $n+1$ and $n(n+2)$ divides $(n+1)^2-n(n+2)=1$, so $n+1$ and $n(n+2)$ are coprime.

Suppose $n(n+1)(n+2)=k^2$. A product of two coprime positive integers is a perfect square only if each factor is a perfect square (each prime appears in only one factor, with an even exponent). So $n(n+2)=n^2+2n$ would be a perfect square. But for $n\ge1$,
\[
n^2<n^2+2n<n^2+2n+1=(n+1)^2,
\]
so $n^2+2n$ lies strictly between two consecutive squares and is not a square. This contradiction shows that $n(n+1)(n+2)$ is never a perfect square. (For $n=0$ the product is $0=0^2$, which is why $n$ must be positive.)
